% Vietnamese translation for OI Public Library
% Source-File: IOI中国国家候选队论文/2003/张宁/张宁.doc
\documentclass[11pt]{article}
\usepackage{oipl-vn}

\title{Nghiên cứu bài toán đoán số}
\author{Zhang Ning}
\date{2003}

\begin{document}
\maketitle

\origtitle{Nghiên cứu bài toán đoán số}
\sourcepath{IOI China candidate team papers / 2003 / Zhang Ning / Zhang Ning.doc}

\section*{Từ khóa}

Suy luận; lồng tư duy; tình huống kết thúc; tình huống loại một; tình
huống loại hai; phân nhóm.

\section*{Tóm tắt}

Trong suy luận lô-gic có một lớp bài toán khá đặc biệt, có thể gọi là bài
toán ``lồng tư duy'': trong đầu \(C\) phải xét \(B\) đang nghĩ như thế nào
về suy nghĩ của \(A\). Những bài toán này thường rất trừu tượng, số trường
hợp cần xét rất nhiều, mạch suy nghĩ cực kỳ phức tạp; nếu phân tích bằng
phương pháp thông thường thì hiệu quả rất kém. Bài viết này lấy một bài toán
điển hình thuộc dạng đó, bài toán đoán số, làm điểm xuất phát. Từ bản chất
của bài toán, ta dùng một cách nhìn mới để phân tích, giải quyết tốt bài toán
và chỉ ra ưu điểm của hướng suy nghĩ mới. Phân tích từ bản chất giúp tránh
được lớp ``lồng tư duy'' bề mặt, đồng thời rút ra nhiều kết luận, làm hiệu
quả giải bài toán tăng rõ rệt. Dưới sự dẫn dắt của cách nghĩ mới, bài toán
gốc được mở rộng sang những tình huống tổng quát hơn. Dù bài toán trở nên
rườm rà và phức tạp hơn, vì nắm được mạch chính của nó, ta vẫn có thể giải
quyết một cách hiệu quả.

\section{Dạng gốc của bài toán}

\subsection*{Mô tả bài toán}

Một giáo sư lô-gic có ba học sinh \(A,B,C\), đều rất giỏi suy luận và tính
nhẩm. Một hôm, giáo sư ra cho ba em một bài toán: ông dán lên trán mỗi người
một tờ giấy và nói với họ rằng trên mỗi tờ đều viết một số nguyên dương, đồng
thời tổng của hai số nào đó bằng số thứ ba. Vì thế, mỗi học sinh nhìn thấy
hai số trên trán hai bạn còn lại, nhưng không nhìn thấy số của mình.

Giáo sư lần lượt hỏi \(A,B,C\): em có đoán được số trên trán mình không? Sau
một số lượt hỏi, khi giáo sư hỏi lại một người, người ấy bỗng mỉm cười tự tin
và nói đúng số ghi trên trán mình.

Bài toán của chúng ta là: chứng minh liệu có nhất định có người đoán được số
của mình hay không; nếu có, hãy tính lượt hỏi sớm nhất mà một người đoán ra,
phân tích toàn bộ quá trình suy luận và tổng kết kết luận thu được.

Trước hết xét một ví dụ đơn giản để quan sát cách mỗi người suy luận. Giả sử
ba số trên trán \(A,B,C\) lần lượt là \(1,2,3\).

Trước tiên hỏi \(A\). Khi đó \(A\) thấy hai số của \(B,C\) là \(2,3\). \(A\)
nhận ra số của mình chỉ có thể là \(3+2=5\), hoặc \(3-2=1\). Không thể biết
đó là \(1\) hay \(5\), nên \(A\) chỉ có thể trả lời ``không''.

Tiếp theo hỏi \(B\). \(B\) nhận ra số của mình chỉ có thể là \(3+1=4\), hoặc
\(3-1=2\). Không biết đó là \(2\) hay \(4\), \(B\) chỉ có thể bắt đầu phân
tích từ câu trả lời của \(A\):

Nếu số của mình là \(2\), thì \(A\) nhìn thấy hai số của \(B,C\) là \(2,3\);
\(A\) sẽ thấy số của mình chỉ có thể là \(3+2=5\) hoặc \(3-2=1\). Không biết
là \(1\) hay \(5\), \(A\) chỉ có thể trả lời ``không''. Điều này trùng với
câu trả lời thực tế của \(A\), không mâu thuẫn, nên \(B\) không thể loại bỏ
trường hợp này.

Nếu số của mình là \(4\), thì \(A\) nhìn thấy hai số của \(B,C\) là \(4,3\);
\(A\) sẽ thấy số của mình chỉ có thể là \(4+3=7\) hoặc \(4-3=1\). Không biết
là \(1\) hay \(7\), \(A\) cũng chỉ có thể trả lời ``không''. Điều này cũng
trùng với câu trả lời thực tế của \(A\), không mâu thuẫn, nên \(B\) không thể
loại bỏ trường hợp này. Vì vậy \(B\) không thể phán đoán, chỉ có thể trả lời
``không''.

Tiếp theo hỏi \(C\). \(C\) nhận ra số của mình chỉ có thể là \(2+1=3\), hoặc
\(2-1=1\). Không biết đó là \(1\) hay \(3\), \(C\) phân tích từ câu trả lời
của \(A\) hoặc \(B\):

Nếu số của mình là \(1\), thì \(A\) sẽ thấy số của mình chỉ có thể là
\(2+1=3\), hoặc \(2-1=1\). Không biết là \(1\) hay \(3\), \(A\) không thể
phán đoán, chỉ có thể trả lời ``không''. Điều này trùng với thực tế, không
mâu thuẫn.

Còn \(B\) sẽ thấy số của mình chỉ có thể là \(1+1=2\), vì số trên trán là số
nguyên dương nên không thể là \(1-1=0\). Do đó \(B\) đáng ra phải trả lời
``có''. Nhưng điều này mâu thuẫn với câu trả lời thực tế của \(B\). \(C\) có
thể nhờ đó loại bỏ trường hợp số của mình là \(1\).

Nếu tiếp tục phân tích trường hợp số của \(C\) là \(3\), ta sẽ không tìm thấy
mâu thuẫn nào, vì đó là tình huống thực tế. Vậy \(C\) phán đoán chính xác số
trên trán mình là \(3\) và trả lời ``có''. Như vậy ở lượt hỏi thứ ba đã có
người đoán được số của mình.

Ta vừa xuất phát từ góc nhìn của từng người để phân tích tình huống các số là
\(1,2,3\). Đây cũng là cách làm phổ biến khi giải những bài suy luận lô-gic
đơn giản. Nhưng nếu quy mô bài toán lớn hơn, độ phức tạp sẽ tăng rất nhanh:
hầu như mỗi lần thêm một tầng suy luận, độ phức tạp lại tăng gấp đôi. Vì
khuôn khổ có hạn, ta không nêu ví dụ phức tạp hơn, nhưng mức độ rườm rà có
thể hình dung được.

Dựa vào loại suy luận rườm rà đó không thể giải tốt bài toán. Nguyên nhân là
có rất nhiều ``lồng tư duy'': trong đầu \(C\) phải xét \(B\) suy nghĩ thế nào
về suy nghĩ của \(A\). Hơn nữa, cách này không rút ra được kết luận có ý
nghĩa phổ quát. Ta hãy đổi cách nghĩ để giải bài toán.

\subsection*{Định nghĩa}

Dưới đây gọi \(A,B,C\) lần lượt là học sinh thứ nhất, thứ hai và thứ ba.

\paragraph{Định nghĩa 1.}
Dùng bộ bốn
\[
  (a_1,a_2,a_3,k),\qquad k\in\{1,2,3\},
\]
để mô tả một tình huống trong quá trình suy luận. Ở đó ba số trên trán lần
lượt là \(a_1,a_2,a_3\), quá trình hỏi bắt đầu từ học sinh thứ nhất, và người
đang được hỏi là học sinh thứ \(k\).

Ký hiệu \(M\) là tập các bộ bốn thỏa điều kiện bài toán:
\[
  M=\{(a_1,a_2,a_3,k)\mid a_i\in\mathbb Z^+,\ k\in\{1,2,3\},
  \text{ một trong }a_1,a_2,a_3\text{ bằng tổng hai số còn lại}\}.
\]

Tính chất \(W\): với bộ bốn \((a_1,a_2,a_3,k)\), học sinh thứ \(k\) có thể
không dựa vào câu trả lời của người khác mà trực tiếp suy ra số của mình. Khi
đó gọi tình huống được mô tả bởi bộ bốn là \term{tình huống kết thúc}. Trong
dạng gốc, nếu \(\{i,j,k\}=\{1,2,3\}\), điều kiện này chính là \(a_i=a_j\).

Tính chất \(S_1\): với bộ bốn \((a_1,a_2,a_3,k)\), có
\[
  a_k=a_i+a_j,\qquad \{i,j,k\}=\{1,2,3\};
\]
gọi tình huống này là \term{tình huống loại một}.

Tính chất \(S_2\): với bộ bốn \((a_1,a_2,a_3,k)\), có
\[
  a_i=a_j+a_k \quad\text{hoặc}\quad a_j=a_i+a_k,
  \qquad \{i,j,k\}=\{1,2,3\};
\]
gọi tình huống này là \term{tình huống loại hai}.

\paragraph{Định nghĩa 2--4.}
Ký hiệu
\[
  M_W=\{x\in M\mid W(x)\},\qquad
  M_1=\{x\in M\mid S_1(x)\},\qquad
  M_2=\{x\in M\mid S_2(x)\}.
\]

\paragraph{Định nghĩa 5.}
Với \(x=(a_1,a_2,a_3,k)\in M\), nếu học sinh thứ \(k\) chưa thể đoán được số
của mình, ghi \(F(x)=\infty\). Nếu học sinh thứ \(k\) có thể đoán được, ghi
\(F(x)\) là tổng số lượt hỏi từ lúc bắt đầu ở học sinh thứ nhất đến lúc số
được đoán ra.

\paragraph{Định nghĩa 6.}
Với bộ bốn \(x=(a_1,a_2,a_3,k)\), ký hiệu
\[
  G(x,R,V)\in\{\mathrm{true},\mathrm{false}\}
\]
để biểu thị rằng khi ba số lần lượt là \(a_1,a_2,a_3\), học sinh thứ \(k\) có
phải là người đầu tiên đoán đúng ở đúng lượt hỏi thứ \(R\), và số trên trán
của người đó có phải là \(V\), hay không.

\paragraph{Định nghĩa 7.}
Với bộ bốn \(x=(a_1,a_2,a_3,k)\), ký hiệu
\[
  H(x,R,V)\in\{\mathrm{true},\mathrm{false}\}
\]
để biểu thị rằng trong tình huống ấy, học sinh thứ \(k\) có thể loại bỏ khả
năng số trên trán mình là \(V\) ở lượt hỏi thứ \(R\), hay không.

\subsection*{Phân tích}

Tiền đề đã biết của bài toán là: tổng của hai số nào đó bằng số thứ ba, và số
trên mỗi tờ giấy đều là số nguyên dương.

Số trên trán mỗi người hoặc bằng tổng hai số còn lại, hoặc bằng hiệu dương
của hai số còn lại. Nếu không thể trực tiếp khẳng định là số nào, chỉ có thể
dùng phép loại trừ. Vì họ sẽ không đoán sai, dưới đây ta chỉ cần xét họ có
thể loại bỏ khả năng khác với tình huống thực tế như thế nào.

Xét bộ bốn \((a_1,a_2,a_3,k)\). Lấy \(i,j\) sao cho
\(\{i,j,k\}=\{1,2,3\}\). Với học sinh thứ \(k\), hai khả năng cho số của
mình là
\[
  a_i+a_j \quad\text{hoặc}\quad |a_i-a_j|.
\]
Khả năng \(a_i+a_j\) không thể bị loại bỏ trực tiếp. Vì số phải dương nên chỉ
khi \(a_i=a_j\) mới có thể trực tiếp loại bỏ khả năng hiệu, tức là đây là
tình huống kết thúc.

Nếu không ở tình huống kết thúc, học sinh thứ \(k\) cần thông qua suy luận để
loại bỏ một trong hai khả năng. Khi học sinh thứ \(k\) giả sử số của mình là
một giá trị nào đó, rồi suy ra đáng lẽ một người khác đã có thể đoán ra trong
hai lượt hỏi trước, thì người đó có thể dùng mâu thuẫn với thực tế ``người
khác chưa đoán ra'' để loại bỏ khả năng ấy. Đây chính là mạch suy luận trong
ví dụ trước đối với \(C\).

Rõ ràng không ai có thể bằng suy luận loại bỏ tình huống thực tế, vì họ không
bao giờ đoán sai. Do đó chỉ cần xét cách loại bỏ khả năng khác với thực tế.
Không mất tính tổng quát, xét tình huống loại một
\[
  a_k=a_i+a_j.
\]
Khi đó học sinh thứ \(k\) cần loại bỏ khả năng
\[
  a_k'=|a_i-a_j|.
\]
Trong các công thức dưới đây, ký hiệu \(\lor\) là phép ``hoặc''.

Với \(k=1\), nếu \(a_2=a_3\) thì kết thúc ở lượt \(1\); nếu \(a_2>a_3\) thì
chuyển sang \((a_2-a_3,a_2,a_3,2)\) sau thêm \(2\) lượt; nếu \(a_3>a_2\) thì
chuyển sang \((a_3-a_2,a_2,a_3,3)\) sau thêm \(1\) lượt:
\[
\begin{aligned}
F(a_1,a_2,a_3,1)&=1 &&(a_2=a_3),\\
F(a_1,a_2,a_3,1)&=2+F(a_2-a_3,a_2,a_3,2) &&(a_2>a_3),\\
F(a_1,a_2,a_3,1)&=1+F(a_3-a_2,a_2,a_3,3) &&(a_3>a_2).
\end{aligned}
\]

Tương tự,
\[
\begin{aligned}
F(a_1,a_2,a_3,2)&=2 &&(a_1=a_3),\\
F(a_1,a_2,a_3,2)&=1+F(a_1,a_1-a_3,a_3,1) &&(a_1>a_3),\\
F(a_1,a_2,a_3,2)&=2+F(a_1,a_3-a_1,a_3,3) &&(a_3>a_1),
\end{aligned}
\]
và
\[
\begin{aligned}
F(a_1,a_2,a_3,3)&=3 &&(a_1=a_2),\\
F(a_1,a_2,a_3,3)&=2+F(a_1,a_2,a_1-a_2,1) &&(a_1>a_2),\\
F(a_1,a_2,a_3,3)&=1+F(a_1,a_2,a_2-a_1,2) &&(a_2>a_1).
\end{aligned}
\]

Như vậy ta có thể chuyển việc xét bộ bốn
\((a_1,a_2,a_3,k)\) sang xét một trong các bộ bốn có số lớn nhất được thay
bằng hiệu của hai số còn lại. Từ đó hình thành một lối suy luận tuyến tính,
không còn là lối suy luận phân nhánh đồ sộ ban đầu. Vì ba số không thể giảm
vô hạn, sau hữu hạn lần biến đổi tất yếu đạt đến tình huống kết thúc.

Kết luận: bất kể ba số thay đổi thế nào, bất kể bắt đầu hỏi từ ai, người có
số lớn nhất trên trán nhất định là người đầu tiên đoán ra số của mình.

Các công thức trên có thể lập trình trực tiếp. Nhờ thiết lập quan hệ truy hồi
tuyến tính, ta tránh được việc quy mô bài toán tăng theo cấp số mũ theo số
lượt hỏi. Cách giải này dựa trên việc phân tích sâu bản chất bài toán. Mạch
suy nghĩ chính có thể tóm tắt là:

\begin{quote}
rút ra tiền đề quan trọng \(\to\) xét tình huống kết thúc \(\to\) lập quan
hệ tương đương cho tình huống chưa kết thúc \(\to\) xét cách quy về tình
huống kết thúc \(\to\) phân trường hợp và chứng minh \(\to\) rút ra kết
luận, viết lại quan hệ tương đương \(\to\) thu được công thức.
\end{quote}
Toàn bộ quá trình bắt đầu từ bản chất của bài toán, thay vì chỉ xuất phát từ
suy nghĩ của từng người. Nhìn từ toàn cục giúp tránh được tầng ``lồng tư
duy'' rườm rà nhất, đồng thời biến quy mô bài toán từ cấp số mũ thành tuyến
tính.

\section{Mở rộng thứ nhất}

\subsection*{Mô tả bài toán}

Một giáo sư lô-gic có \(n\ge 3\) học sinh rất giỏi suy luận và tính nhẩm.
Ông dán lên trán mỗi người một số nguyên dương và nói rằng một trong các số
bằng tổng của \(n-1\) số còn lại. Mỗi học sinh nhìn thấy \(n-1\) số của những
người khác, nhưng không nhìn thấy số của mình.

Giáo sư lần lượt hỏi \(n\) học sinh: em có đoán được số của mình không? Sau
một số lượt hỏi, một người khi được hỏi lại đã đoán đúng số của mình.

Bài toán là chứng minh liệu có nhất định có người đoán ra hay không; nếu có,
tính lượt hỏi sớm nhất, phân tích quá trình suy luận và tổng kết kết luận.
Ta sửa các định nghĩa của trường hợp \(n=3\) như sau.

\subsection*{Định nghĩa}

Dùng bộ \((n+1)\)-phần tử
\[
  (a_1,a_2,\ldots,a_n,k)
\]
để mô tả một tình huống, trong đó \(a_i\) là số trên trán học sinh thứ \(i\),
quá trình hỏi bắt đầu từ học sinh thứ nhất, và người đang được hỏi là học
sinh thứ \(k\).

Gọi \(M_n\) là tập các bộ thỏa
\[
  a_i\in\mathbb Z^+,\qquad
  \exists p:\ a_p=\sum_{i\ne p}a_i.
\]
Tình huống kết thúc \(W\) là tình huống học sinh đang được hỏi có thể trực
tiếp xác định số của mình. Tình huống loại một \(S_1\) là
\[
  a_k=\sum_{i\ne k}a_i,
\]
còn tình huống loại hai \(S_2\) là tình huống số lớn nhất nằm ở một học sinh
khác. Các tập \(M_W,M_1,M_2\), hàm \(F\), các mệnh đề \(G,H\) được định nghĩa
tương tự như trong dạng gốc.

Với \(A=(a_1,\ldots,a_n)\), đặt
\[
  S_k=\sum_{i\ne k}a_i.
\]
Với học sinh thứ \(k\), số của mình có hai khả năng:
\[
  S_k \quad\text{hoặc}\quad
  2a_v-S_k,
\]
trong đó \(v\ne k\) là chỉ số của số lớn nhất trong các số mà học sinh thứ
\(k\) nhìn thấy. Nếu \(2a_v\le S_k\), khả năng thứ hai không dương và bị loại
trực tiếp; khi đó đây là tình huống kết thúc.

Nếu \(2a_v>S_k\), để loại bỏ khả năng khác với thực tế, ta xét bộ mới nhận
được bằng cách thay số của học sinh thứ \(k\) bởi
\[
  a_k'=2a_v-S_k.
\]
Tức là
\[
  (a_1,\ldots,a_k,\ldots,a_n,k)
  \longrightarrow
  (a_1,\ldots,2a_v-S_k,\ldots,a_n,v).
\]
Độ lệch lượt hỏi từ \(v\) tới \(k\) là
\[
  d(v,k)=
  \begin{cases}
    k-v, & v<k,\\
    n-v+k, & v>k.
  \end{cases}
\]

Vì vậy, nếu \(k\) là chỉ số của số lớn nhất trong \(A\), \(v\) là chỉ số của
số lớn thứ hai và \(S=\sum_{i\ne k}a_i\), truy hồi tính số lượt hỏi là
\[
F(A,k)=
\begin{cases}
  k, & 2a_v\le S,\\
  d(v,k)+F(A',v), & 2a_v>S,
\end{cases}
\]
trong đó \(A'\) được tạo từ \(A\) bằng cách thay
\[
  a_k\leftarrow 2a_v-S.
\]
Do ta chỉ cần xét trường hợp số lớn nhất là số bằng tổng các số còn lại, chỉ
số \(k\) được xác định trực tiếp từ \(A\), nên \(F(A,k)\) có thể viết gọn
thành \(F(A)\).

Kết luận: bất kể \(n\) số thay đổi thế nào, bất kể bắt đầu hỏi từ ai, người
có số lớn nhất trên trán nhất định là người đầu tiên đoán ra. Sử dụng công
thức trên có thể lập trình giải bài toán. Trường hợp \(n=3\) cung cấp một đối
chiếu rất tốt, giúp ta giải trường hợp tổng quát tương đối nhẹ nhàng; bản
thân điều này cũng dựa trên phân tích của trường hợp \(n=3\).

\section{Mở rộng thứ hai}

\subsection*{Mô tả bài toán}

Một giáo sư lô-gic có \(n\ge 3\) học sinh. Ông dán lên trán mỗi người một số
nguyên dương, rồi chia họ thành hai nhóm. Một nhóm có \(m\) người,
\[
  1\le m\le n-1,
\]
và các học sinh không biết cách chia nhóm. Giáo sư cho biết tổng các số trên
trán hai nhóm bằng nhau. Mỗi học sinh nhìn thấy \(n-1\) số của các bạn còn
lại, nhưng không nhìn thấy số của mình.

Giáo sư lần lượt hỏi \(n\) học sinh: em có đoán được số của mình không? Sau
một số lượt hỏi, một người đoán đúng số trên trán mình. Bài toán vẫn là
chứng minh có nhất định có người đoán ra hay không; nếu có, tính lượt hỏi sớm
nhất, phân tích quá trình suy luận và rút ra kết luận.

Khi \(n=3\), \(m\) chỉ có thể là \(2\), tức là dạng gốc. Khi \(m=n-1\), ta
thu được mở rộng thứ nhất. Vì vậy phần dưới chỉ xét \(n>3\) và \(m<n-1\).

Trong bài toán này, khi một người phán đoán số của mình, giá trị có thể phụ
thuộc vào cách chia nhóm. Với học sinh thứ \(k\), người đó nhìn thấy tất cả
trừ số của mình. Nếu giả sử một cách chia nhóm, từ điều kiện ``tổng hai nhóm
bằng nhau'' có thể tính ra số của mình. Có
\[
  \binom{n-1}{m}\quad\text{hoặc}\quad \binom{n-1}{m-1}
\]
cách chọn nhóm liên quan đến người thứ \(k\), tùy việc nhóm của \(k\) có
kích thước \(m\) hay \(n-m\). Do đó có nhiều giá trị khả dĩ tương ứng; trong
đó có thể có giá trị trùng nhau hoặc không dương.

\subsection*{Định nghĩa}

Ta vẫn dùng bộ
\[
  (a_1,a_2,\ldots,a_n,k)
\]
để mô tả tình huống, và dùng các tính chất \(W,S_1,S_2\), các tập
\(M_W,M_1,M_2\), hàm \(F\), mệnh đề \(G,H\) theo nghĩa tương tự.

Với \(i\ne k\), đặt
\[
  A_{\max}=\max_i a_i,\qquad
  I_{\max}=\{i\mid a_i=A_{\max}\}.
\]
Với một cách chia nhóm giả định \(X\) của học sinh thứ \(k\), ký hiệu
\[
  G_k(X)
\]
là hiệu ``tổng nhóm đối diện trừ tổng nhóm cùng phía, không kể số chưa biết
của học sinh thứ \(k\)''; giá trị khả dĩ của số trên trán học sinh thứ \(k\)
được ký hiệu là
\[
  V_k(X).
\]
Nếu \(V_k(X)\le 0\), cách chia nhóm ấy bị loại trực tiếp.

Sắp các số mà học sinh thứ \(k\) nhìn thấy theo thứ tự giảm:
\[
  A_{t_1}\ge A_{t_2}\ge \cdots\ge A_{t_{n-1}}.
\]
Xét cách chia đặc biệt \(T\): chọn \(m\) học sinh có số lớn nhất trong
\(\{t_1,\ldots,t_{n-1}\}\) làm một nhóm, nhóm còn lại chứa học sinh thứ
\(k\). Khi đó giá trị suy ra dưới cách chia \(T\) là
\[
  V_k(T)=\sum_{i=1}^m A_{t_i}
        -\sum_{i=m+1}^{n-1}A_{t_i}.
\]
Nếu \(V_k(T)\le 0\), tình huống tương ứng bị loại trực tiếp. Nếu
\[
  0<V_k(T)<a_k,
\]
thì đây là một khả năng khác với thực tế và cần được loại bỏ bằng suy luận.

\subsection*{Phân tích các tình huống chính}

Đối với một cách chia bất kỳ \(X\), giả sử nhóm khác với học sinh thứ \(k\)
gồm \(m\) người có chỉ số \(u_1,\ldots,u_m\). Khi đó
\[
  V_k(X)=\sum_{r=1}^m a_{u_r}
  -\sum_{i\notin\{k,u_1,\ldots,u_m\}} a_i.
\]
Nếu nhóm cùng phía với học sinh thứ \(k\) có \(m\) người còn lại, công thức
tương ứng là
\[
  V_k(X)=
  \sum_{i\notin\{k,u_1,\ldots,u_m\}} a_i
  -\sum_{r=1}^m a_{u_r}.
\]
Hai công thức chỉ khác nhau ở việc nhóm nào được coi là nhóm đối diện với
\(k\).

Tình huống có thể kết thúc trực tiếp khi các giá trị khả dĩ dương của
\(V_k(X)\) chỉ còn một giá trị. Nếu có nhiều giá trị dương khác nhau, học
sinh thứ \(k\) phải thông qua câu trả lời của những người trước đó để loại bỏ
từng giá trị không đúng. Về bản chất, mỗi lần loại bỏ một giá trị \(V\) lại
tạo ra một tình huống mới trong đó số của học sinh thứ \(k\) được thay bởi
\(V\), rồi xét xem trong tình huống mới ấy có người khác đáng lẽ đã đoán ra
sớm hơn hay không.

Với cách chia đặc biệt \(T\), nếu giá trị \(V_k(T)\) là một khả năng cần loại
bỏ, ta chuyển sang tình huống
\[
  A'=(a_1,\ldots,a_{k-1},V_k(T),a_{k+1},\ldots,a_n).
\]
Nếu trong \(A'\) có một học sinh khác \(p\) đáng lẽ đoán ra trước lượt hỏi
của \(k\), thì \(k\) loại bỏ được khả năng \(V_k(T)\). Ngược lại, nếu không
tạo ra mâu thuẫn, \(k\) chưa thể đoán ra.

Khi \(m=n-1\), phép biến đổi này suy biến đúng về mở rộng thứ nhất: chỉ có
hai khả năng bản chất, và quá trình suy luận tuyến tính. Khi \(m<n-1\), số
cách chia tăng lên; các khả năng không còn tạo thành một đường suy luận duy
nhất mà tạo thành một cây suy luận.

Xét riêng các trường hợp có nhiều học sinh cùng mang số lớn nhất.
\begin{itemize}
  \item Nếu tất cả \(n\) số đều bằng nhau, tình huống kết thúc ngay.
  \item Nếu có nhiều hơn hai số lớn nhất, hoặc \(m>n/2\), các giá trị khả dĩ
        không đủ để tạo ra mâu thuẫn duy nhất; không ai có thể là người đầu
        tiên đoán ra trong nhánh ấy.
  \item Nếu đúng hai số lớn nhất và các số còn lại không bằng nhau, nhánh suy
        luận cũng không quy về một tình huống kết thúc duy nhất.
  \item Nếu đúng hai số lớn nhất, các số còn lại bằng nhau, \(n=4\), và hai
        số còn lại nhỏ hơn một nửa số lớn nhất theo điều kiện cần, ta nhận
        được một tình huống hai đôi bằng nhau. Tình huống này được gọi là
        \term{loại A}; qua mỗi bước suy luận, một trong bốn số giảm, nên sau
        hữu hạn bước sẽ đạt tình huống kết thúc.
\end{itemize}

Đối với trường hợp chỉ có một số lớn nhất, giả sử chỉ số của nó là \(q\). Bỏ
qua \(a_q\) và sắp \(n-1\) số còn lại giảm dần:
\[
  b_1\ge b_2\ge\cdots\ge b_{n-1}.
\]
Điều kiện cần để có thể tiếp tục suy luận là
\[
  a_q=\sum_{i=1}^{m}b_i-\sum_{i=m+1}^{n-1}b_i.
\]
Nếu không thỏa điều kiện này, không có cách chia thực tế phù hợp với giả định
đang xét. Nếu thỏa, với mỗi tập \(C\) gồm \(m\) người, đặt
\[
  s(C)=\sum_{i\in C}a_i.
\]
Giá trị cần xét cho học sinh thứ \(q\) là
\[
  V_q(C)=2s(C)-\sum_{i\ne q}a_i.
\]
Chỉ những \(C\) thỏa
\[
  0<V_q(C)<a_q
\]
mới tạo ra khả năng khác với thực tế cần loại bỏ. Với mỗi khả năng như vậy,
ta đệ quy trên tình huống thay \(a_q\) bởi \(V_q(C)\). Nếu mọi khả năng sai
đều dẫn tới mâu thuẫn trước lượt hỏi của \(q\), thì \(q\) có thể đoán ra;
nếu tồn tại một khả năng không loại bỏ được, thì \(q\) chưa thể là người đầu
tiên đoán ra.

Từ các thảo luận trên, tiêu chuẩn tiếp tục suy luận có thể viết như sau:
với học sinh thứ \(k\), cần có \(V_k(T)<a_k\), tức tồn tại giá trị khả dĩ
dương nhỏ hơn giá trị thực; đồng thời với mọi cách chia khả dĩ \(C\) tạo ra
giá trị \(0<V_k(C)<a_k\), tình huống đệ quy tương ứng phải cho phép một
người khác đoán ra trước lượt hỏi hiện tại. Khi đó học sinh thứ \(k\) mới có
thể loại bỏ toàn bộ khả năng sai.

Trong trường hợp \(n=4,m=2\), nếu rơi vào loại A nói trên, nhất định có người
đoán được. Nếu không rơi vào loại A, các giá trị khả dĩ còn lại vẫn giảm qua
mỗi bước suy luận; sau hữu hạn bước sẽ đạt tình huống kết thúc. Vì thế mọi
tình huống của \(n=4,m=2\) đều có người đoán ra.

Đối với mở rộng thứ nhất, tức \(m=n-1\), cũng luôn có người đoán ra. Nhưng
với mở rộng thứ hai nói chung, quá trình suy luận có thể phân nhánh nhiều
lần; nó không còn là suy luận tuyến tính mà trở thành một cấu trúc cây. Do
các cách chia nhóm rất nhiều và tiêu chuẩn phán đoán cũng khá phức tạp, việc
tính giá trị \(F\) bằng tay không còn khả thi, nhưng đã có thể giải bằng lập
trình: ghi nhớ các tình huống đã xét, liệt kê các cách chia khả dĩ, loại bỏ
các giá trị không dương, rồi đệ quy trên mọi giá trị sai cần loại bỏ.

\section*{Kết luận}

Bài viết đã phân tích sâu một bài toán suy luận lô-gic. Từ góc nhìn toàn
cục, ta xét mối liên hệ bản chất của bài toán, thay vì chỉ xuất phát từ suy
nghĩ riêng của từng người. Cách làm này đơn giản hóa phần ``lồng tư duy''
rườm rà nhất, đồng thời thiết lập được quan hệ truy hồi, tránh cho quy mô bài
toán tăng nhanh theo số tầng suy luận và giải quyết bài toán hiệu quả. Dựa
trên đối chiếu đó, bài toán được mở rộng sang những tình huống tổng quát hơn.
Đặc biệt trong mở rộng thứ hai, phần thảo luận rất rườm rà, nhưng tư tưởng
cốt lõi vẫn nhất quán. Đây là một phương pháp có thể tham khảo khi giải các
bài toán suy luận lô-gic.

\section*{Tài liệu tham khảo}

\begin{itemize}
  \item \emph{Phân tích CTSC 2001}.
  \item Fu Qingxiang, Wang Xiaodong, \emph{Thuật toán và cấu trúc dữ liệu}.
\end{itemize}

\end{document}
